% ============================================================================
%  RADS — MDPI resubmission skeleton
%  Target journal: Journal of Cybersecurity and Privacy (JCP).  To switch,
%  change the first option in \documentclass below to e.g. futureinternet,
%  information, applsci, electronics.
%
%  HOW TO USE THIS FILE
%  --------------------
%  1. Open the OFFICIAL MDPI Overleaf template for your target journal.
%     (Overleaf: "MDPI Article".) It ships Definitions/mdpi.cls and the
%     bibliography style — this skeleton assumes those files are present.
%  2. Replace that template's main.tex body with this one, keeping its
%     Definitions/ folder. Put references.bib alongside.
%  3. If a macro name below differs from your template version, reconcile
%     with the template's own example file — MDPI updates the class
%     periodically and I can't verify the exact current macro set from here.
%  4. TODO markers show where content still has to be written.
% ============================================================================

\documentclass[jcp,article,submit,pdftex,moreauthors]{Definitions/mdpi}

% amsmath/amssymb are loaded by the MDPI class; if compiling standalone, add:
% \usepackage{amsmath,amssymb}

% ---------------------------------------------------------------- front matter
\Title{Taming Ransomware Decisions: A Prospect-Theoretic Decision-Support Model that Counteracts Attacker Framing}
\TitleCitation{Taming Ransomware Decisions: A Prospect-Theoretic Decision-Support Model that Counteracts Attacker Framing}

\newcommand{\orcidauthorA}{0000-0000-0000-0000} % TODO: your ORCID

\Author{Pranjal Sharma $^{1,}$*\orcidA{}}
\AuthorNames{Pranjal Sharma}
\AuthorCitation{Sharma, P.}

\address{%
$^{1}$ \quad Department of Computer Science, Heidelberg University, Heidelberg, Germany}

\corres{Correspondence: pranjaldub@gmail.com}

\abstract{% TODO: rewrite once results exist. Target structure for MDPI:
% (1) problem + urgency, (2) the gap: existing pay/no-pay guidance ignores that
% attackers weaponise cognitive bias, (3) what we do: formalise the attack as
% a Prospect-Theory reference shift + framing distortion and define RADS, a
% decision-support model that reverses that distortion, (4) how we evaluate it
% (simulation of N scenarios vs an expected-value optimum + retrospective cases
% + sensitivity analysis), (5) headline quantitative result (regret reduction),
% (6) one-line implication.
Ransomware operators do not only encrypt data; they engineer the victim's
decision. This paper formalises that manipulation using Prospect Theory---
modelling the ransom demand as an attacker-induced shift of the decision
reference point together with an urgency-driven distortion of outcome
probabilities---and introduces \emph{Ransomware Risk Analysis and Decision
Support} (RADS), a model that reverses the induced distortion and recommends
whether paying is defensible. We evaluate RADS against a risk-neutral
expected-value optimum on a simulated population of incidents, on documented
real-world cases, and under sensitivity analysis of its parameters.
% TODO: state the quantitative outcome.
}

\keyword{ransomware; prospect theory; behavioural economics; security decision support;
expected value; framing effect; ISO/IEC 27001; DORA}

\begin{document}
\maketitle

% ============================================================================
\section{Introduction}
% Keep the motivating incidents (MGM/Caesars, share-price drop, the 17\% rise)
% but end the intro with an explicit, testable contribution list.
\cite{bib_statista}\cite{bib_cvedetails}\cite{bib_techcrunch}

\textbf{Contributions.} % <- reviewers look for this.
\begin{enumerate}
  \item We formalise the attacker's manipulation as two Prospect-Theory
        operators: a reference-point shift and an urgency-driven probability
        distortion (Section~\ref{sec:method}).
  \item We define RADS, which inverts those operators and compares the
        de-biased valuation against a risk-neutral expected-value benchmark.
  \item We evaluate RADS by simulation, retrospective case analysis, and
        sensitivity analysis (Sections~\ref{sec:eval}--\ref{sec:results}),
        rather than a single illustrative example.
\end{enumerate}

% ============================================================================
\section{Related Work}\label{sec:related}
% NEW SECTION — the single biggest structural gap in the arXiv version.
% Cover three strands and position RADS against each:
%   (a) economics / game theory of ransom payment \cite{bib_cartwright2019,bib_laszka2017};
%   (b) behavioural economics of security decisions (Prospect Theory) \cite{bib_prospecttheory,bib_cpt1992};
%   (c) incident-response standards as decision baselines \cite{bib_iso,bib_dora}.
% Gap statement: prior work models the *economics* of paying, but not the
% attacker's deliberate exploitation of cognitive bias, nor a mechanism to
% reverse it. That gap is RADS.
% TODO: 15--30 references. Verify every reconstructed citation in references.bib.

% ============================================================================
\section{Background}\label{sec:background}

\subsection{Incident-response baseline (ISO/IEC 27001:2022)}
% Condensed from the original §1.1: A.16 (incident management) and A.17
% (continuity) establish that a responsible owner and a restoration plan must
% already exist. Keep short — this is a prerequisite, not the contribution.
\cite{bib_iso}

\subsection{Prospect Theory}
Prospect Theory \cite{bib_prospecttheory} models choice through a value
function defined over gains and losses relative to a reference point, and a
probability-weighting function \cite{bib_cpt1992}. The value function is
S-shaped and loss-averse:
\begin{equation}\label{eq:value}
v(x)=
\begin{cases}
x^{\alpha}, & x \ge 0,\\[4pt]
-\lambda\,(-x)^{\beta}, & x < 0,
\end{cases}
\end{equation}
with $\alpha,\beta \approx 0.88$ and loss-aversion $\lambda \approx 2.25$
\cite{bib_cpt1992}. Objective probabilities are transformed by
\begin{equation}\label{eq:weight}
w(p)=\frac{p^{\gamma}}{\left(p^{\gamma}+(1-p)^{\gamma}\right)^{1/\gamma}},
\end{equation}
with $\gamma \approx 0.61$, so that small probabilities are over-weighted and
large ones under-weighted. Equations~\eqref{eq:value}--\eqref{eq:weight} are
the machinery the original draft named but did not use; they are central here.

% ============================================================================
\section{Method: The RADS Model}\label{sec:method}

\subsection{Decision setting}
For each affected application $a$ the organisation chooses
$k\in\{\mathrm{pay},\mathrm{rec}\}$ (pay the ransom, or self-recover). Both
outcomes are uncertain. Let $r_0$ denote the organisation's true pre-attack
state, used as the neutral reference point.

\subsection{Expected-value benchmark (ground truth)}
Self-recovery has expected monetary loss
\begin{equation}\label{eq:erec}
\mathbb{E}\!\left[\ell_a^{\mathrm{rec}}\right]=
C^{\mathrm{down}}_a\,t^{\mathrm{rec}}_a
+C^{\mathrm{lab}}_a
+(1-\phi_a)\,V^{\mathrm{data}}_a
+p^{\mathrm{leak}}_{\mathrm{np}}\,D^{\mathrm{leak}}_a,
\end{equation}
where $\phi_a\in[0,1]$ is backup coverage (the ``feasibility'' factor),
$C^{\mathrm{down}}_a$ the downtime cost rate, $t^{\mathrm{rec}}_a$ the recovery
time, $V^{\mathrm{data}}_a$ the value of unrecoverable current data,
$D^{\mathrm{leak}}_a$ the leak damage, and $p^{\mathrm{leak}}_{\mathrm{np}}$ the
probability of leak given non-payment. Paying has expected loss
\begin{equation}\label{eq:epay}
\mathbb{E}\!\left[\ell_a^{\mathrm{pay}}\right]=
R_a
+p_{\mathrm{key}}\,C^{\mathrm{down}}_a t^{\mathrm{pay}}_a
+(1-p_{\mathrm{key}})\,\mathbb{E}\!\left[\ell_a^{\mathrm{rec}}\right]
+p_{\mathrm{re}}\,C^{\mathrm{re}}_a
+p^{\mathrm{leak}}_{\mathrm{p}}\,D^{\mathrm{leak}}_a,
\end{equation}
where $R_a$ is the ransom, $p_{\mathrm{key}}$ the probability the decryptor
actually works, $p_{\mathrm{re}}$ the re-extortion probability, and
$C^{\mathrm{re}}_a$ its cost. Note that Eq.~\eqref{eq:epay} makes explicit what
attacker framing hides: paying does not guarantee recovery. The risk-neutral
optimum is
\begin{equation}\label{eq:opt}
d_a^{\star}=\arg\min_{k\in\{\mathrm{pay},\mathrm{rec}\}}\mathbb{E}\!\left[\ell_a^{(k)}\right].
\end{equation}

\subsection{Perceived value and the attacker's manipulation}
A decision-maker does not evaluate Eqs.~\eqref{eq:erec}--\eqref{eq:epay}
directly; they evaluate them through Prospect Theory:
\begin{equation}\label{eq:ptval}
V_a^{\mathrm{PT}}\!\left(k \mid r,\{p_i\}\right)=\sum_i w(p_i)\,v\!\left(x_i^{(k)}-r\right),
\end{equation}
where $x_i^{(k)}$ are the monetary outcomes of action $k$. The attacker applies
two operators, parameterised by an urgency level $u\in[0,1]$ set by the
deadline/doubling scheme:
\begin{equation}\label{eq:manip}
\tilde r = r_0-\Delta_{\mathrm{ref}}(u),\qquad
\tilde p_{\mathrm{key}} = p_{\mathrm{key}}+u\,(1-p_{\mathrm{key}}),\qquad
\tilde p_{\mathrm{re}} = (1-u)\,p_{\mathrm{re}}.
\end{equation}
The reference shift $\tilde r$ reframes catastrophe as the baseline (Figure~2 of
the original draft), so paying reads as a \emph{gain}; the urgency terms
over-weight the chance the key works and suppress re-extortion. The manipulated
choice is $d_a^{\mathrm{PT}}=\arg\max_k V_a^{\mathrm{PT}}(k \mid \tilde r,\{\tilde p_i\})$.

\subsection{The RADS correction}
RADS reverses Eq.~\eqref{eq:manip}: it restores the reference to $r_0$ using
prior-incident evidence (the ``reference effect'') and replaces manipulated
probabilities with empirical priors, then recommends
\begin{equation}\label{eq:rads}
d_a^{\mathrm{RADS}}=\arg\max_k V_a^{\mathrm{PT}}\!\left(k \mid r_0,\{p_i\}\right).
\end{equation}
RADS therefore removes the \emph{attacker-induced} distortion while remaining
consistent with Eq.~\eqref{eq:opt}; the two benchmarks $d_a^{\star}$ and
$d_a^{\mathrm{RADS}}$ are reported side by side.

\subsection{Portfolio prioritisation by criticality}
% This preserves the original criticality idea, on a single consistent scale
% kappa_a in (0,1]. (The arXiv version silently switched the scale from
% {1,0.5,0.25} to {5,2}; that is fixed here.)
A screening score aggregates the normalised factors $V_f$ with weights
$\omega_f$:
\begin{equation}\label{eq:df}
\mathrm{DF}_a=\sum_{f}\omega_f\,\rho_f\,V_f,\qquad \sum_f\omega_f=1,
\end{equation}
where $\rho_f\in\{+1,-1\}$ fixes each factor's direction (e.g. higher backup
coverage must \emph{lower} the incentive to pay --- a sign error in the original
formula). Applications are then ranked by a criticality-adjusted score,
\begin{equation}\label{eq:adj}
\widehat{\mathrm{DF}}_a=\mathrm{DF}_a\left(1+\frac{\kappa_a-\bar\kappa}{\sum_j \kappa_j}\right),
\end{equation}
with $\kappa_a$ the application criticality and $\bar\kappa$ its mean.

% ============================================================================
\section{Evaluation Design}\label{sec:eval}
% This section replaces the single "Crown Financing GmbH" example.

\subsection{Simulation study}
We sample a population of $N$ incident scenarios, drawing
$R_a, p_{\mathrm{key}}, p_{\mathrm{re}}, p^{\mathrm{leak}}, \phi_a, C^{\mathrm{down}}_a,
t^{\mathrm{rec}}_a, V^{\mathrm{data}}_a$ from ranges anchored in public incident
data \cite{bib_coveware,bib_sophos,bib_ibmcost,bib_enisa}. For each scenario we
compute three decisions --- the risk-neutral optimum $d^{\star}$
(Eq.~\eqref{eq:opt}), the manipulated Prospect-Theory decision $d^{\mathrm{PT}}$,
and the RADS decision $d^{\mathrm{RADS}}$ (Eq.~\eqref{eq:rads}) --- and report
decision agreement, over-pay and under-pay rates, and expected regret.

\subsection{Regret metric}
\begin{equation}\label{eq:regret}
\mathrm{Regret}(d)=\mathbb{E}\!\left[\ell_a^{(d)}\right]-\min_k \mathbb{E}\!\left[\ell_a^{(k)}\right].
\end{equation}
Headline claim to test: $\overline{\mathrm{Regret}}(d^{\mathrm{RADS}})\ll
\overline{\mathrm{Regret}}(d^{\mathrm{PT}})$ across the population.

\subsection{Sensitivity analysis}
Monte-Carlo perturbation of the weights $\omega_f$ and PT parameters
$(\alpha,\beta,\lambda,\gamma)$; report the fraction of scenarios whose decision
is stable under $\pm x\%$ perturbation and a tornado plot of factor influence.
This turns the ``arbitrary weights/threshold'' criticism into a characterised
result.

\subsection{Retrospective cases}
Apply RADS to 8--15 documented incidents with known outcomes (e.g. MGM,
Caesars, Colonial Pipeline, Norsk Hydro, Maersk/NotPetya, Baltimore, Atlanta);
compare the recommendation to what each organisation did. Presented as
illustrative, not statistical --- state the small-$N$, confounded-outcome
limitation explicitly.

% ============================================================================
\section{Results}\label{sec:results}
% TODO: figures + tables produced by the simulation notebook.
% - Fig: regret distribution, naive / PT-biased / RADS.
% - Table: decision-agreement and over/under-pay rates.
% - Fig: sensitivity tornado + stability curve.
% - Table: retrospective cases (org | did | RADS says | post-hoc better call).

% ============================================================================
\section{Discussion}\label{sec:discussion}
% - Legality/sanctions (OFAC-type exposure), insurance, law-enforcement
%   reporting: a model that can output "pay" must engage these seriously.
% - The ethical stance: RADS de-biases; it does not endorse funding crime.
% - Applicability to DORA (Art. 8 asset identification, Art. 29 concentration
%   risk): the same PT risk scoring drives a compliance-oriented asset register.
% - Limitations: parameter elicitation, single-organisation reference data.

% ============================================================================
\section{Conclusions}\label{sec:conclusion}
% Tighten the original conclusion; restate the tested claim and future work
% (human-subjects framing study; the DORA registry application).

% ---------------------------------------------------------------- back matter
\funding{This research received no external funding.}

\acknowledgments{The author thanks Prof.\ Dr.\ Rima-Maria Rahal, Department of
Psychology, Heidelberg University, for formal feedback.}

\conflictsofinterest{The author declares no conflict of interest.}

\begin{adjustwidth}{-\extralength}{0cm}
\reftitle{References}
% MDPI templates typically use: \externalbibliography{yes}
\bibliography{references}
\end{adjustwidth}

\end{document}
